% !TeX root = ../main.tex
\section{Introduction}
\label{sec:intro}

Folk games are part of intangible cultural heritage, which the UNESCO
Convention describes as \emph{transmitted from generation to
generation}~\citep{unesco2003ich}. Many Vietnamese folk board games have no
written rules; the rules are kept only in the memory of the people who
played them. When those players grow up and stop playing, a game can vanish
without a trace. Some games have been recorded
and digitised, such as Shoulder-Pole Chess (our translation of its Vietnamese name) from
Quang Nam~\citep{mytour_coganh}. Many
others have not.

\textbf{Tactical Formation} belongs to the second group. It is a board game
played by the children of Dong Bua hamlet, Tam Quan commune, Tam Dao
district, Vinh Phuc province; \emph{Tactical Formation} is the author's
English rendering of its village name (Section~\ref{sec:provenance}). The
author played the game throughout grades 5 to 9. Today almost nobody in the village remembers
the game, and the author had partly forgotten it too. This paper comes from
the wish to keep Tactical Formation from being forgotten: to write its rules down, to
model it mathematically, and to put it into software that anyone can play.
The author did \emph{not} create Tactical Formation and does not know who did.
Section~\ref{sec:provenance} describes its origin and how it was recorded.

In terms of mechanics, every move in Tactical Formation is a \emph{jump} over an
adjacent piece, including one's own. The jumped-over piece stays on the
board. From the landing square, the player \emph{sandwiches} the line of
enemy pieces lying between the moved piece and a friendly piece, and that
whole line changes side. Every move therefore has two effects at once: the
new position of the piece shapes future jumps, and the sandwich formed from
that position flips pieces immediately. The player must weigh both. Taken
one at a time, the jump resembles Halma and the capture resembles Reversi.
Section~\ref{sec:related} compares Tactical Formation with these games in detail.

The paper makes the following contributions:
\begin{enumerate}[label=(\roman*)]
  \item To our knowledge, the first complete and consistent written record
        of the rules of Tactical Formation for 2--4 players, together with a record of
        the game's origin. The 3-player starting configuration is recorded
        exactly as in the folk game, even though it is unbalanced; see
        Section~\ref{sec:limitations}.
  \item A description of the two-layer board structure (vertices and
        centres), with the movement and capture directions of each kind of piece.
  \item A formal mathematical model for the 2-player mode, together with
        the 28-piece conservation invariant (\emph{capture = change of
        side}) and upper bounds on the number of actions and states.
  \item A companion AI based on Minimax and Paranoid Minimax, with the
        reasons for not using MCTS.
  \item An empirical evaluation of the game and the AI via branching
        factor and game length (2--4 players), games between bots (2-player
        mode), and the bots' think time.
  \item An open-source software release, as a form of digitisation and
        preservation, so that the community can play and study the game.
\end{enumerate}

The remainder of the paper is organised as follows.
Section~\ref{sec:provenance} describes the origin of the game and the
recording method. Section~\ref{sec:related} compares Tactical Formation with related
games. Section~\ref{sec:rules} gives a complete specification of the rules
through board images. Section~\ref{sec:variants} describes the multiplayer
modes. Section~\ref{sec:ai} presents the companion-AI algorithm.
Section~\ref{sec:impl} describes the software. Section~\ref{sec:eval}
gives the empirical evaluation. Section~\ref{sec:limitations} states the
limitations, and Section~\ref{sec:conclusion} concludes. The full formal
mathematical model for the 2-player mode is presented in
Appendix~\ref{app:formal}.
